% Include as a fragment in the patterns handbook; the root supplies the
% KACTL package, Prelude, landscape page, 10pt type, and two-column layout.
\section{Model and pattern cases}

These cases train the step before coding: translate the story into objects,
states, and constraints, then choose the smallest tool that preserves what
the answer needs. The first two cases are real contest problems with links to
their statements and official solution material. The last two are synthetic
models designed to expose reusable reductions.

\subsection{Case: intervals - Regional 2024 J, \emph{Jigsaw of Shadows}}
\label{case:intervals}

The statement gives an angle $\theta$, then $N$ flatlanders at distinct road
positions $X_i$, each of height $H_i$. Each casts a shadow eastward; report
the total road length covered, counting overlap once. The input bound is
$N\le 10^5$, $10\le\theta\le80$ degrees, $X_i\le3\cdot10^5$,
$H_i\le1000$, and output error may be $10^{-4}$. See the
\href{https://codeforces.com/gym/105505/attachments/download/28244/contest_onesided.pdf}{official
statement} and \href{https://codeforces.com/gym/105505/attachments/download/28286/icpcla24_solution_sketches.pdf}{official
solution sketches}.

The story looks like geometry, but the requested object is one-dimensional
length. A right triangle gives shadow length $L_i=H_i\cot(\theta)$, so each
person contributes the closed interval $[X_i,X_i+L_i]$. Once represented this
way, none of the other people interact geometrically: only interval overlap
matters. This is a useful recognition cue for rays, projected objects,
coverage windows, and occupied time ranges: derive each range, then ask
whether the answer is the union, intersection, count, or maximum depth.

Convert degrees before calling trigonometry:
\verb@rad = theta * acos(-1.0) / 180.0@ and
\verb@L = H * cos(rad) / sin(rad)@. Build the vector of endpoint pairs and
call \verb@finals_core::merge_intervals<double>(intervals)@. That helper
sorts by left then right endpoint, skips zero-length intervals, and merges
overlap or touching endpoints. Sum \verb@r-l@ over the returned disjoint
intervals. The exact helper returns intervals rather than their measure, so
it can also support later queries about gaps or merged components.

Synthetic trace with $\theta=60^\circ$:
$(X,H)=(50,150),(0,100),(100,200)$. Their lengths are about $86.6025$,
$57.7350$, and $115.4701$. After sorting, the ranges are
$[0,57.7350]$, $[50,136.6025]$, and $[100,215.4701]$. Each overlaps the
current covered prefix, so the merged range is $[0,215.4701]$ and the answer
is about $215.47$. Summing raw shadow lengths counts the shared road more
than once.

Sorting dominates at $O(N\log N)$ time; merging and summing use $O(N)$ time
and memory. A sweep over compressed real endpoints can also work, but each compressed
gap must retain its original length. Counting compressed indices would
lose geometry: integer road positions do not imply integer shadow endpoints. A sweep over start/end events also works, but adds
event ordering and active-count logic when a single sorted merge is enough.
If the task changes to count how many shadows cover each query point, retain
the event sweep or a sorted endpoint structure. If it asks for the uncovered
gaps inside a finite road, clip first and keep the merged ranges. The sample
allows floating output; keep the computation in double and do not round
individual endpoints before merging.

Practice: two shadows are $[2,7]$ and $[7,9]$. What is the union length?
Answer: $7$, because the touching point contributes no length twice; the
helper merges them to $[2,9]$.

\subsection{Case: flip - Masters Mexico 2022 F, \emph{Flip}}
\label{case:flip}

The real problem starts from an all-white $3\times3$ board. Clicking a cell
toggles it and its orthogonal neighbors that exist; find the minimum clicks
for a given target. The statement guarantees a solution and has at most 50
boards. Read the \href{https://acm.iteso.mx/share/ICPCMXMasters2022_problems.pdf}{official
problem book}, page 9 (problem F). This is a compact instance of a broader
recognition rule: if an action is binary, its effects are fixed, actions
commute, and the number of actions is tiny, enumerate action subsets before
reaching for search or linear algebra.

Number cells row-major from 0 to 8. Represent a board as a 9-bit unsigned
mask, with bit $r\cdot3+c$ set when that cell is black. For each cell $i$,
precompute an effect mask containing $i$ and its in-bounds up/down/left/right
neighbors. A press subset $p$ changes state by XOR with the XOR of all
selected effect masks. Since each click toggles and clicks commute, pressing
one cell twice cancels; therefore some minimum solution uses each cell zero
or one times. Every candidate is one of $2^9=512$ masks. Count its set bits
and keep the least count among candidates whose combined effect equals the
target.

The implementation module provides
\verb@finals_cases::min_toggle_count(effects, start, target)@. Here
\verb@effects[i]@ is the mask toggled by action $i$; \verb@start@ and
\verb@target@ use the same bit layout. It returns the minimum number of
actions or $-1$ if unreachable, with a maximum of 20 actions. For this
problem use start mask zero. Parse each row by shifting left or setting bit
$3r+c$ for each \verb@*@, and construct effects using bounds checks rather
than wrapping neighbors across a row.
\sourceimport{finals-patterns}{ToggleSubsets.h}

Synthetic trace:
\begin{lstlisting}
*..
**.
*..
\end{lstlisting}
The target mask has bits $0,1,3$ set. Pressing cell 0 toggles exactly cells
$0,1,3$, so this subset has size 1 and reaches the target. No zero-click
subset can produce a nonzero target, so the minimum is 1. The next sample
board can be tested by the same loop; the routine does not rely on visual
symmetry or board rotation.

The cost is $O(2^k\cdot k)$ time and $O(k)$ effect storage for $k$ actions;
the stated helper limit is $k\le20$. Exhaustive enumeration is clearer and
more reliable here than a generic graph shortest path, because the cost is
the number of chosen independent actions and each subset has one direct
result. Gaussian elimination over $\mathbb F_2$ can find a solution for a
larger linear system, but it does not by itself minimize the number of set
variables; minimum-weight solutions require extra work. If the board grows
to a narrow strip, a row-mask DP may exploit width. If actions have different
costs or dependencies, subset enumeration can still work at small $k$, but
the objective or legal-subset test must change. Do not use a bit-shift mask
when $k$ exceeds the width of its integer type; use multiple words or a
different representation.

Practice: two actions toggle masks $0011$ and $0110$, and the target is
$0101$ from zero. Is it reachable? Yes. Pressing both actions produces
$0011\mathbin{\mathtt{XOR}}0110=0101$, so the minimum is two. Effects combine
by XOR, not OR.

\subsection{Case: matching or flow - assigning work with quotas}
\label{case:matching}

Synthetic example: jobs $A$ and $B$ each need one worker. Worker Red can do
either job and worker Blue can do only $A$; each worker can take at most one
job. Can both jobs be assigned? Model jobs and individual workers as the two
sides of a bipartite graph; an edge means the assignment is allowed. A
matching selects edges with no shared endpoint. Here $A\to$ Blue and
$B\to$ Red is a complete assignment.

For a standard one-to-one case, build adjacency list \verb@g[job]@ of worker
indices, initialize the right-match array to $-1$, and call
\verb@hopcroftKarp(g, rightMatch)@. Its returned size is the largest number
of jobs that can be assigned. It runs in $O(E\sqrt V)$, where $E$ is the
number of allowed pairs and $V$ the number of vertices; each job and worker
must be used at most once, and the objective is unweighted cardinality.
Use this direct helper when those constraints exactly match the story.

Now extend the example: Red is a team with two workers, both able to do $A$
or $B$; Blue still has one worker who can do only $A$. A team node with
capacity one would incorrectly restrict Red to one job. Model source $s$,
job nodes, team nodes, and sink $t$. Add $s\to$ each job with capacity 1;
each allowed job-to-team edge has capacity 1; add Red-to-$t$ with capacity 2
and Blue-to-$t$ with capacity 1. A unit of integral flow along
$s\to job\to team\to t$ is one assignment. The invariant is flow
conservation at every job/team node, plus edge capacities; integer
capacities guarantee an integral maximum flow. Run Dinic from $s$ to $t$.
Both source edges saturated means both jobs were assigned.

Send $B$ to Red and $A$ to Blue, for flow value 2 and a complete assignment.
If Blue is removed, Red's two workers can still take both jobs. This shows
why a team quota belongs on the team-to-sink edge: source-to-job edges model
job demand, while team-to-sink edges model available worker capacity.

Use the supplied \verb@Dinic f(n); f.addEdge(u,v,cap); f.calc(s,t)@ helper
when capacities or shared bottlenecks matter. Its flow cost depends on the
network and capacities; for the unit-capacity bipartite construction the
specialized Hopcroft--Karp bound is $O(E\sqrt V)$ and is usually the simpler
guarantee. If preferences have costs, use min-cost flow rather than plain
maximum flow. If each person may receive several jobs, increase that worker-to-sink
capacity; if a team quota has a lower bound, ordinary max flow alone does not
enforce it. A weighted matching helper may fit a dense assignment with
one-to-one capacities, but the flow model is more adaptable when quotas or
intermediate resource limits appear.

Practice: remove worker Blue from the two-job example. What is the maximum
assignment size? Answer: 2. Both jobs can go to Red's two workers.

\subsection{Case: capacity - binary search plus a greedy proof}
\label{case:capacity}

Synthetic scheduling task: jobs arrive in a fixed order with nonnegative
durations. Split them into at most $k$ nonempty contiguous groups, and
minimize the largest group sum. Jobs cannot be reordered or split. Example
durations are $[4,2,3,5]$ with $k=2$. First define the decision question: for
capacity $C$, can
the jobs be assigned to at most $k$ groups so every group sum is at most
$C$? Increasing $C$ can only make a previously feasible grouping remain
feasible, so the predicate has the form false, false, ..., true, true.

For a fixed $C$, scan left to right and append each job to the current group
if it fits; otherwise close that group and start the next with the job. This
greedy scan uses the fewest possible groups: its first group consumes the
longest prefix whose sum is at most $C$. Any feasible partition's first
group ends no later; remove the greedy prefix and repeat the same argument
on the suffix. Thus if greedy needs at most $k$ groups, a valid partition
exists; if it needs more, no partition can do better. The invariant during
the scan is that the current group is the maximal feasible prefix not yet
closed and the group count is the minimum needed for the processed prefix.

Use the largest job as the lower bound: any feasible capacity must be at
least this large, though the bound itself need not be infeasible. The sum of
all jobs is a feasible upper bound (one group always fits). Call
\verb@finals_cases::min_contiguous_capacity(a, groups)@; it assumes
nonnegative jobs, a sum fitting in \verb@long long@, and returns 0 for an
empty list. The helper finds the minimum feasible integer capacity by
monotone binary search and the greedy test. With $n$ jobs and total sum $S$,
this is $O(n\log(S+1))$ time and $O(1)$ extra space, aside from the input.

Trace the example. At $C=7$, greedy forms $[4,2]$, then $[3]$, then $[5]$:
three groups, so infeasible. At $C=8$, it forms $[4,2]$ and $[3,5]$: two
groups, so feasible. Binary search narrows the integer boundary between
these cases; 8 is the minimum. A tempting wrong predicate is ``greedy uses
exactly $k$ groups.'' It can become true and later false as capacity grows,
so it is not the monotone decision question. For $[2,2,2]$ and $k=2$, it
uses three groups at $C=2$, exactly two at $C=4$, and one at $C=6$; this
false/true/false sequence cannot be binary-searched as a first-true test.

The local helper avoids a common midpoint bug by using
\verb@lo + (hi-lo)/2@. Its half-open search discards tested infeasible
capacities; the original upper bound is a proven feasible fallback. The
lower bound itself need not be infeasible. The exact API returns
the minimum capacity, so the caller does not need to reproduce the search.
If groups must be exactly $k$, at-most-$k$ feasibility needs an extra proof:
with nonnegative jobs, any group can be split until there are $k$ groups
provided $k\le n$ and empty groups are forbidden; this follows by splitting
multi-job groups without increasing their sums. If negative durations are
allowed, this greedy prefix proof fails: extending a group can reduce its
sum, so taking the longest feasible prefix need not minimize the group count.
Feasibility itself remains monotone in $C$: every partition feasible at $C$
remains feasible at a larger capacity. If jobs may be
reordered, sorting or a different partition model may be needed. If the
objective is minimizing the largest number of jobs rather than total weight,
the feasibility test changes even though binary search may still apply.

Practice: durations $[6,1,4]$, at most two groups. Is capacity 5 feasible?
Answer: no; the first job alone exceeds 5.
Capacity 6 is feasible with $[6]$, $[1,4]$, so the minimum is 6.
\sourceimport{finals-patterns}{CapacitySearch.h}
